\documentclass[11pt]{amsart}

\usepackage{amsmath,amssymb,amsthm,mathtools}
\usepackage{enumitem}

\newtheorem{theorem}{Theorem}[section]
\newtheorem{corollary}[theorem]{Corollary}
\newtheorem{lemma}[theorem]{Lemma}
\newtheorem{proposition}[theorem]{Proposition}
\newtheorem{definition}[theorem]{Definition}
\newtheorem{remark}[theorem]{Remark}

\DeclareMathOperator{\oc}{oc}
\DeclareMathOperator{\lcp}{lcp}
\DeclareMathOperator{\dimH}{dim_H}

\begin{document}

\title{Metric and Extremal Theory of Online Previous-Copy Compression}

\author{Luca Blanchi}

\date{June 2026}

\begin{abstract}
Let $b\geq2$ be an integer alphabet size. We study the extremal and metric behaviour of the online previous-copy parsing complexity $\oc(W)$ of a finite word $W\in\{0,\ldots,b-1\}^N$. In this model a phrase is either a literal symbol or an exact copy whose source starts earlier in the word; self-overlap is allowed.

We prove the sharp universal upper bound
\[
\oc(W)\leq (1+o(1))\frac{N}{\log_b N}
\]
uniformly for all words $W$ of length $N$, and show that this is best possible:
\[
\max_{|W|=N}\oc(W)
=
(1+o(1))\frac{N}{\log_b N}.
\]
For Lebesgue-almost every $\alpha\in[0,1]$, if $w_N(\alpha)$ denotes the prefix of length $N$ of the base-$b$ expansion of $\alpha$, then
\[
\oc(w_N(\alpha))
=
(1+o(1))\frac{N}{\log_b N}.
\]

We also determine the finite-word entropy of low online previous-copy complexity. For $0\leq \kappa\leq1$, let
\[
\mathcal C_N(\kappa)
=
\left\{
W\in\{0,\ldots,b-1\}^N:
\oc(W)\leq \kappa\frac{N}{\log_b N}
\right\}.
\]
Then, for $0<\kappa\leq1$,
\[
\lim_{N\to\infty}
\frac1N\log_b|\mathcal C_N(\kappa)|
=
\kappa.
\]
Consequently, the Hausdorff spectrum is exact:
\[
\dimH
\left\{
\alpha\in[0,1]:
\liminf_{N\to\infty}
\frac{\oc(w_N(\alpha))\log_b N}{N}
\leq\kappa
\right\}
=
\kappa
\qquad(0\leq\kappa\leq1).
\]

Finally, we record parallel extremal and almost-sure estimates for the normalized substring complexity
\[
\delta(W)=\max_{1\leq k\leq |W|}\frac{p_W(k)}{k}
\]
and for the minimum string attractor size $\gamma(W)$:
\[
\delta(w_N(\alpha))
=
\gamma(w_N(\alpha))
=
(1+o(1))\frac{N}{\log_b N}
\]
for almost every $\alpha$, and the same asymptotic scale holds in the worst case.
\end{abstract}

\maketitle

\section{Introduction}

The online previous-copy parsing model is a structural abstraction of left-to-right dictionary compression. A finite word is parsed from left to right into phrases. A phrase is either a literal symbol or a copy of an earlier substring. The source is required to start earlier than the target, but it may overlap the target. This convention includes the self-referential behaviour familiar from variants of LZ77.

A companion paper studied this model on digit prefixes of algebraic irrational numbers and proved the lower bound
\[
\oc(w_N(\alpha))=\omega(\log N)
\]
for every algebraic irrational $\alpha$. The present note is independent of the Diophantine part. Its purpose is to determine the natural extremal and metric scale of $\oc(W)$ for arbitrary and typical finite words.

The main scale is
\[
\frac{N}{\log_b N}.
\]
This scale appears in three complementary ways.

First, every word of length $N$ has an online previous-copy parsing with at most
\[
(1+o(1))\frac{N}{\log_b N}
\]
phrases. Second, by a counting argument, this is sharp in the worst case. Third, the same asymptotic holds for almost every base-$b$ expansion.

The proof of the universal upper bound is elementary. Choose a block length
\[
L=\lfloor \log_b N-2\log_b\log_b N\rfloor.
\]
Parsing greedily, any phrase shorter than $L$, except possibly near the end of the word, must start at a previously unseen block of length $L$. There are at most $b^L$ such blocks, while all other phrases have length at least $L$. This gives
\[
\oc(W)\leq \frac{N}{L}+b^L+L+O(1).
\]

The lower bounds come from counting parse descriptions. The number of words of length $N$ admitting an online previous-copy parse with $z$ phrases is at most
\[
z\left(\frac{eN}{z}\right)^z(2bN)^z.
\]
For
\[
z\sim \kappa \frac{N}{\log_b N},
\]
this is $b^{(\kappa+o(1))N}$. This gives both the worst-case lower bound and the entropy/Hausdorff spectrum.

The last part treats two standard repetitiveness measures: the normalized substring complexity
\[
\delta(W)=\max_k p_W(k)/k
\]
and the minimum string attractor size $\gamma(W)$. The inequalities
\[
\delta(W)\leq\gamma(W)\leq\oc(W)
\]
connect these measures to the online previous-copy model. A standard collision estimate for random words gives the matching lower bound for $\delta$, and hence for $\gamma$, almost surely.

Throughout the paper, logarithms without subscript are natural, while $\log_b$ denotes logarithm in base $b$.

\section{Online previous-copy parsings}

Let
\[
W=W[1]W[2]\cdots W[N]
\]
be a finite word over the alphabet
\[
\{0,\ldots,b-1\}.
\]

\begin{definition}
An online previous-copy parsing of $W$ is a factorization
\[
W=F_1F_2\cdots F_z
\]
with boundaries
\[
0=n_0<n_1<\cdots<n_z=N,
\]
where
\[
F_j=W[n_{j-1}+1,n_j].
\]
Each phrase is of one of the following two types.

First, $F_j$ may be a literal, in which case
\[
|F_j|=1.
\]

Second, $F_j$ may be a copy. Writing
\[
s=n_{j-1},
\qquad
 t=n_j,
\qquad
\ell=t-s,
\]
this means that there exists a source position $p$ with
\[
1\leq p\leq s
\]
such that
\[
W[p+h]=W[s+1+h]
\qquad
(0\leq h<\ell).
\]
The source may overlap the target.
\end{definition}

Let
\[
\oc(W)
\]
be the minimum number of phrases in an online previous-copy parsing of $W$.

\begin{remark}
The first phrase must be a literal. The definition is structural: source positions and lengths are not charged as bits. The quantity $\oc(W)$ is therefore a phrase-complexity measure, not a bit-complexity measure.
\end{remark}

\section{A universal upper bound}

We first prove that every word has an online previous-copy parsing with about $N/\log_b N$ phrases.

\begin{theorem}[Universal upper bound]
Uniformly for all words $W\in\{0,\ldots,b-1\}^N$,
\[
\oc(W)
\leq
(1+o(1))\frac{N}{\log_b N}.
\]
More precisely, for every integer $L\geq1$,
\[
\oc(W)
\leq
\frac{N}{L}+b^L+L+O(1),
\]
and with
\[
L=\left\lfloor \log_b N-2\log_b\log_b N\right\rfloor
\]
this gives the displayed asymptotic.
\end{theorem}

\begin{proof}
Fix $L$. We construct a parsing by the following greedy rule. Suppose the current position is $i$. If $i+L-1\leq N$ and the block
\[
W[i,i+L-1]
\]
has an occurrence starting at some position $p<i$, we make a copied phrase of length $L$, with source $p$. Otherwise we make a literal phrase of length $1$. Once fewer than $L$ symbols remain, we parse the rest as literals.

Call a phrase short if it is a literal produced before the final $L$ positions. If a short phrase starts at position $i\leq N-L+1$, then the block
\[
W[i,i+L-1]
\]
has no earlier occurrence. Hence two different short phrase starts $i<j\leq N-L+1$ give two different length-$L$ blocks. Otherwise, if
\[
W[i,i+L-1]=W[j,j+L-1],
\]
then at position $j$ the parser could have copied at least $L$ symbols from source $i$, contradicting that the phrase at $j$ was short.

There are at most $b^L$ distinct blocks of length $L$. Hence the number of short phrases before the final $L$ positions is at most $b^L$. The final part contributes at most $L$ literal phrases.

All copied phrases produced by the algorithm have length $L$, so there are at most $N/L$ such phrases. Therefore
\[
\oc(W)
\leq
\frac{N}{L}+b^L+L+O(1).
\]

Now take
\[
L=\left\lfloor \log_b N-2\log_b\log_b N\right\rfloor.
\]
Then
\[
b^L\leq \frac{N}{(\log_b N)^2},
\]
and
\[
\frac{N}{L}
=
(1+o(1))\frac{N}{\log_b N}.
\]
Also
\[
b^L+L=o\left(\frac{N}{\log_b N}\right).
\]
Thus
\[
\oc(W)
\leq
(1+o(1))\frac{N}{\log_b N}.
\]
\end{proof}

\section{Counting words with short parsings}

Let
\[
\mathcal W_{N,z}
=
\left\{
W\in\{0,\ldots,b-1\}^N:
\oc(W)\leq z
\right\}.
\]

\begin{lemma}[Counting parse descriptions]
For $1\leq z\leq N/2$,
\[
|\mathcal W_{N,z}|
\leq
z\left(\frac{eN}{z}\right)^z(2bN)^z.
\]
Consequently, if
\[
z_N=\left\lfloor \kappa\frac{N}{\log_b N}\right\rfloor
\]
with fixed $0<\kappa\leq1$, then
\[
|\mathcal W_{N,z_N}|
\leq
b^{(\kappa+o(1))N}.
\]
\end{lemma}

\begin{proof}
Fix $m\leq z$. We count a superset of words admitting a parsing with $m$ phrases.

The boundaries are determined by the $m-1$ internal cut positions, so there are
\[
\binom{N-1}{m-1}
\]
choices.

For each phrase, choose whether it is a literal or a copy, giving at most $2^m$ choices. Assign to every phrase a symbol in the alphabet, even if the phrase is a copy; this gives at most $b^m$ choices. Assign also to every phrase a source position in $\{1,\ldots,N\}$, even if the phrase is a literal; this gives at most $N^m$ choices.

Thus the number of descriptions with $m$ phrases is at most
\[
\binom{N-1}{m-1}(2bN)^m.
\]
Every such description determines at most one word. Indeed, literals prescribe their symbols, and copied phrases are deterministic once the already parsed prefix is known. If a source overlaps the target, the copied symbols are still determined progressively because the source starts earlier than the target.

Therefore
\[
|\mathcal W_{N,z}|
\leq
\sum_{m\leq z}
\binom{N-1}{m-1}(2bN)^m.
\]
For $m\leq z\leq N/2$,
\[
\binom{N-1}{m-1}
\leq
\left(\frac{eN}{m}\right)^m.
\]
The function
\[
m\mapsto
\left(\frac{eN}{m}\right)^m(2bN)^m
\]
is increasing for $1\leq m\leq N/2$. Hence
\[
|\mathcal W_{N,z}|
\leq
z\left(\frac{eN}{z}\right)^z(2bN)^z.
\]

Now take
\[
z_N=\left\lfloor \kappa\frac{N}{\log_b N}\right\rfloor.
\]
Taking logarithms in base $b$,
\[
\log_b |\mathcal W_{N,z_N}|
\leq
\log_b z_N
+
z_N\log_b\frac{eN}{z_N}
+
z_N\log_b(2bN).
\]
The last term is
\[
z_N\log_b N+O(z_N)
=
(\kappa+o(1))N.
\]
The middle term is
\[
z_N\log_b\frac{eN}{z_N}
=
O\left(\frac{N}{\log N}\log\log N\right)
=
o(N).
\]
Also $\log_b z_N=o(N)$. Therefore
\[
\log_b |\mathcal W_{N,z_N}|
\leq
(\kappa+o(1))N,
\]
as claimed.
\end{proof}

\section{Worst-case complexity}

\begin{theorem}[Worst-case asymptotics]
\[
\max_{W\in\{0,\ldots,b-1\}^N}\oc(W)
=
(1+o(1))\frac{N}{\log_b N}.
\]
\end{theorem}

\begin{proof}
The upper bound follows from the universal upper bound.

For the lower bound, fix $0<\kappa<1$, and set
\[
z_N=\left\lfloor \kappa\frac{N}{\log_b N}\right\rfloor.
\]
By the counting lemma,
\[
|\mathcal W_{N,z_N}|
\leq
b^{(\kappa+o(1))N}.
\]
Since the total number of words of length $N$ is $b^N$, for all large $N$ there exists a word $W$ with
\[
\oc(W)>z_N.
\]
Thus
\[
\max_{|W|=N}\oc(W)
\geq
\kappa\frac{N}{\log_b N}(1-o(1)).
\]
Letting $\kappa\uparrow1$ gives the lower bound.
\end{proof}

\section{Metric law for online previous-copy complexity}

For $\alpha\in[0,1]$, let
\[
w_N(\alpha)
\]
be the prefix of length $N$ of the base-$b$ expansion of $\alpha$. We ignore the countable set of numbers with two base-$b$ expansions.

\begin{theorem}[Almost-sure asymptotic]
For Lebesgue-almost every $\alpha\in[0,1]$,
\[
\oc(w_N(\alpha))
=
(1+o(1))\frac{N}{\log_b N}.
\]
Equivalently,
\[
\lim_{N\to\infty}
\frac{\oc(w_N(\alpha))\log_b N}{N}
=
1
\]
for almost every $\alpha$.
\end{theorem}

\begin{proof}
The upper bound holds for every word by the universal upper bound.

For the lower bound, fix $0<\kappa<1$, and set
\[
z_N=\left\lfloor \kappa\frac{N}{\log_b N}\right\rfloor.
\]
The set of $\alpha$ such that
\[
\oc(w_N(\alpha))\leq z_N
\]
is a union of base-$b$ cylinders of length $N$, one for each word in $\mathcal W_{N,z_N}$. Each cylinder has Lebesgue measure $b^{-N}$. Hence its measure is at most
\[
b^{-N}|\mathcal W_{N,z_N}|
\leq
b^{-(1-\kappa+o(1))N}.
\]
This is summable in $N$. By the Borel--Cantelli lemma, for almost every $\alpha$, the event
\[
\oc(w_N(\alpha))\leq \kappa\frac{N}{\log_b N}
\]
occurs for only finitely many $N$. Therefore
\[
\liminf_{N\to\infty}
\frac{\oc(w_N(\alpha))\log_b N}{N}
\geq \kappa
\]
for every $0<\kappa<1$, and hence the liminf is at least $1$. The limsup is at most $1$ by the universal upper bound.
\end{proof}

\section{Finite entropy of low-complexity words}

For $0<\kappa\leq1$, define
\[
\mathcal C_N(\kappa)
=
\left\{
W\in\{0,\ldots,b-1\}^N:
\oc(W)\leq
\kappa\frac{N}{\log_b N}
\right\}.
\]

\begin{theorem}[Entropy spectrum]
For $0<\kappa\leq1$,
\[
\lim_{N\to\infty}
\frac1N\log_b|\mathcal C_N(\kappa)|
=
\kappa.
\]
For $\kappa\geq1$, the limit is $1$.
\end{theorem}

\begin{proof}
The upper bound for $0<\kappa\leq1$ follows directly from the counting lemma:
\[
|\mathcal C_N(\kappa)|
\leq
b^{(\kappa+o(1))N}.
\]

For the lower bound, fix $0<\eta<\kappa$, and put
\[
M=\lfloor(\kappa-\eta)N\rfloor.
\]
For every word
\[
U\in\{0,\ldots,b-1\}^M,
\]
construct a word $W(U)$ of length $N$ as follows. First write $U$. Then continue periodically with period $U$ until the word has total length $N$.

Different $U$'s give different $W(U)$'s, because the first $M$ symbols of $W(U)$ are $U$. Hence this construction gives $b^M$ distinct words.

By the universal upper bound applied to $U$,
\[
\oc(U)\leq
(1+o(1))\frac{M}{\log_b M}
\]
uniformly in $U$. After $U$ has been parsed, the periodic continuation is obtained with one copied phrase, using source position $1$ and allowing self-overlap. Thus
\[
\oc(W(U))
\leq
(1+o(1))\frac{M}{\log_b M}+1.
\]
Since
\[
M=(\kappa-\eta+o(1))N
\qquad\text{and}\qquad
\log_b M\sim\log_b N,
\]
we get
\[
\oc(W(U))
\leq
(\kappa-\eta+o(1))\frac{N}{\log_b N}
\leq
\kappa\frac{N}{\log_b N}
\]
for all sufficiently large $N$. Hence
\[
|\mathcal C_N(\kappa)|
\geq
b^M
=
b^{(\kappa-\eta+o(1))N}.
\]
Letting $\eta\to0$ gives
\[
\liminf_{N\to\infty}
\frac1N\log_b|\mathcal C_N(\kappa)|
\geq \kappa.
\]

For $\kappa\geq1$, the universal upper bound implies that all words are eventually counted up to $o(1)$ in the threshold, and the entropy limit is $1$.
\end{proof}

\section{Hausdorff spectrum}

For $0\leq\kappa\leq1$, define
\[
F_\kappa
=
\left\{
\alpha\in[0,1]:
\liminf_{N\to\infty}
\frac{\oc(w_N(\alpha))\log_b N}{N}
\leq \kappa
\right\}.
\]

\begin{theorem}[Hausdorff spectrum]
For $0\leq\kappa\leq1$,
\[
\dimH F_\kappa=\kappa.
\]
For $\kappa\geq1$, $F_\kappa$ has full Lebesgue measure and hence Hausdorff dimension $1$.
\end{theorem}

\begin{proof}
We first prove the upper bound. Fix $\eta>0$. If $\alpha\in F_\kappa$, then for infinitely many $N$,
\[
\oc(w_N(\alpha))
\leq
(\kappa+\eta)\frac{N}{\log_b N}.
\]
Thus $F_\kappa$ is contained in the limsup of the union of cylinders of length $N$ corresponding to words in
\[
\mathcal C_N(\kappa+\eta).
\]
By the entropy upper bound,
\[
|\mathcal C_N(\kappa+\eta)|
\leq
b^{(\kappa+\eta+o(1))N}.
\]
Each cylinder has diameter comparable to $b^{-N}$. If $s>\kappa+\eta$, then
\[
\sum_N
|\mathcal C_N(\kappa+\eta)| b^{-sN}
<
\infty.
\]
Hence the $s$-dimensional Hausdorff measure of $F_\kappa$ is zero. Therefore
\[
\dimH F_\kappa\leq \kappa+\eta.
\]
Letting $\eta\to0$ gives
\[
\dimH F_\kappa\leq\kappa.
\]

For the lower bound, assume $0<\kappa\leq1$. We construct a Cantor set contained in $F_\kappa$. Choose a sequence of integers $M_j\to\infty$ growing so rapidly that the total length constructed before stage $j$ is $o(M_j/\log M_j)$.

Suppose that before stage $j$ a prefix of length $R_{j-1}$ has been constructed. Choose freely a word
\[
U_j\in\{0,\ldots,b-1\}^{M_j}.
\]
After writing $U_j$, continue periodically with period $U_j$ until the total length reaches
\[
N_j
=
R_{j-1}
+
\left\lfloor\frac{M_j}{\kappa}\right\rfloor.
\]
Since $R_{j-1}=o(M_j)$,
\[
N_j\sim \frac{M_j}{\kappa}.
\]

At time $N_j$, the already existing prefix before $U_j$ contributes $o(M_j/\log M_j)$ phrases. The block $U_j$ can be parsed using the universal upper bound:
\[
\oc(U_j)
\leq
(1+o(1))\frac{M_j}{\log_b M_j}.
\]
The periodic continuation after $U_j$ is copied with one phrase, since self-overlap is allowed. Thus
\[
\oc(w_{N_j})
\leq
(1+o(1))\frac{M_j}{\log_b M_j}.
\]
Because
\[
N_j\sim\frac{M_j}{\kappa}
\qquad\text{and}\qquad
\log_b N_j\sim\log_b M_j,
\]
we obtain
\[
\frac{\oc(w_{N_j})\log_b N_j}{N_j}
\leq
\kappa+o(1).
\]
Hence every point in the constructed Cantor set lies in $F_\kappa$.

It remains to estimate its dimension. At stage $j$, the number of independent choices is $b^{M_j}$. Since the sequence $M_j$ grows very fast, the contribution of all previous stages is negligible compared with $M_j$. The cylinders at stage $j$ have length $b^{-N_j}$. A standard mass distribution argument gives
\[
\dimH
\geq
\liminf_{j\to\infty}
\frac{\sum_{i\leq j}M_i}{N_j}
=
\liminf_{j\to\infty}
\frac{M_j+o(M_j)}{M_j/\kappa}
=
\kappa.
\]
Thus
\[
\dimH F_\kappa\geq\kappa.
\]
The case $\kappa=0$ follows from the upper bound and the fact that $F_0$ is nonempty, for example it contains eventually periodic expansions. Therefore $\dimH F_0=0$.
\end{proof}

\section{Normalized substring complexity and string attractors}

For a finite word $W$, let $p_W(k)$ be the number of distinct factors of $W$ of length $k$. Define
\[
\delta(W)=\max_{1\leq k\leq |W|}\frac{p_W(k)}{k}.
\]

Let $\gamma(W)$ denote the size of the smallest string attractor of $W$. Recall that a set of positions
\[
\Gamma\subseteq\{1,\ldots,|W|\}
\]
is a string attractor if every distinct factor of $W$ has an occurrence crossing at least one position of $\Gamma$.

\begin{lemma}
For every finite word $W$,
\[
\delta(W)\leq \gamma(W)\leq \oc(W).
\]
\end{lemma}

\begin{proof}
First let $\Gamma$ be a string attractor for $W$. Fix $k$. Every distinct length-$k$ factor has an occurrence crossing some position of $\Gamma$. A fixed position can be crossed by at most $k$ length-$k$ factors. Therefore
\[
p_W(k)\leq |\Gamma|k.
\]
Taking the minimum over $\Gamma$ and then the maximum over $k$ gives
\[
\delta(W)\leq \gamma(W).
\]

Now take an online previous-copy parsing of $W$ with phrase starts
\[
n_0+1,n_1+1,\ldots,n_{z-1}+1.
\]
Let $\Gamma$ be this set of phrase starts. We claim that $\Gamma$ is a string attractor.

Let $U$ be any factor of $W$, and choose its leftmost occurrence. If this occurrence does not cross a phrase start, it lies strictly inside a single phrase. If that phrase is a literal, this is impossible unless $|U|=1$, in which case the occurrence crosses the phrase start. If the phrase is copied, the source of the copy gives an earlier occurrence of $U$, contradicting the choice of the leftmost occurrence. Hence every factor has an occurrence crossing a phrase start, and $\Gamma$ is an attractor. Thus
\[
\gamma(W)\leq z.
\]
Taking the minimum over parsings gives
\[
\gamma(W)\leq \oc(W).
\]
\end{proof}

\subsection{Universal and worst-case bounds}

\begin{lemma}[Universal upper bound for $\delta$]
Uniformly for all words $W$ of length $N$,
\[
\delta(W)
\leq
(1+o(1))\frac{N}{\log_b N}.
\]
\end{lemma}

\begin{proof}
For every $k$,
\[
p_W(k)\leq \min\{b^k,N-k+1\}\leq \min\{b^k,N\}.
\]
Thus
\[
\frac{p_W(k)}{k}
\leq
\frac{\min\{b^k,N\}}{k}.
\]
If $k\leq \log_b N$, then the maximum of $b^k/k$ in this range occurs at $k=\lfloor\log_b N\rfloor+O(1)$, and is
\[
(1+o(1))\frac{N}{\log_b N}.
\]
If $k\geq \log_b N$, then
\[
\frac{N}{k}\leq \frac{N}{\log_b N}.
\]
Taking the maximum over $k$ proves the claim.
\end{proof}

\begin{theorem}[Worst-case for $\delta$ and $\gamma$]
\[
\max_{|W|=N}\delta(W)
=
(1+o(1))\frac{N}{\log_b N},
\]
and
\[
\max_{|W|=N}\gamma(W)
=
(1+o(1))\frac{N}{\log_b N}.
\]
\end{theorem}

\begin{proof}
The upper bound for $\delta$ is the universal bound just proved. The upper bound for $\gamma$ follows from
\[
\gamma(W)\leq \oc(W)
\]
and the universal upper bound for $\oc$.

For the lower bounds, it is enough to show that there exist words $W$ with
\[
\delta(W)\geq(1-o(1))\frac{N}{\log_b N}.
\]
This follows, for example, from the almost-sure lower bound for $\delta$ proved in the next subsection. Since
\[
\delta(W)\leq\gamma(W),
\]
the same examples give the lower bound for $\gamma$.
\end{proof}

\subsection{Almost-sure behaviour}

\begin{theorem}[Almost-sure behaviour of $\delta$ and $\gamma$]
For Lebesgue-almost every $\alpha\in[0,1]$,
\[
\delta(w_N(\alpha))
=
(1+o(1))\frac{N}{\log_b N},
\]
and
\[
\gamma(w_N(\alpha))
=
(1+o(1))\frac{N}{\log_b N}.
\]
\end{theorem}

\begin{proof}
The upper bound for $\delta$ is universal. The upper bound for $\gamma$ follows from
\[
\gamma(W)\leq\oc(W)
\]
and the almost-sure asymptotic for $\oc$.

It remains to prove the almost-sure lower bound for $\delta$. Fix $\varepsilon>0$, and set
\[
k_N=\left\lceil(1+\varepsilon)\log_b N\right\rceil.
\]
For a random word $W_N$ of length $N$, let $C_N$ be the number of pairs $1\leq i<j\leq N-k_N+1$ such that
\[
W_N[i,i+k_N-1]=W_N[j,j+k_N-1].
\]
For any pair $i<j$, the probability of equality is at most $b^{-k_N}$, even when the two blocks overlap. Indeed, if $j-i=h<k_N$, the equality forces a word of length $k_N+h$ to have period $h$, and the probability is at most $b^{-k_N}$. Therefore
\[
\mathbb E C_N
\ll
N^2b^{-k_N}
\ll
N^{1-\varepsilon}.
\]

Choose a subsequence
\[
N_m=\lfloor m^q\rfloor
\]
with $q\varepsilon>2$. By Markov's inequality,
\[
\mathbb P(C_{N_m}>\eta N_m)
\leq
\frac{\mathbb E C_{N_m}}{\eta N_m}
\ll
N_m^{-\varepsilon}
\ll
m^{-q\varepsilon}.
\]
The last series is summable. By Borel--Cantelli,
\[
C_{N_m}=o(N_m)
\]
almost surely.

The number of distinct length-$k_{N_m}$ factors in the prefix of length $N_m$ is at least
\[
N_m-k_{N_m}+1-C_{N_m}
=
(1-o(1))N_m,
\]
because the number of repeated occurrences is bounded by the number of colliding pairs.

Now let $N_m\leq N<N_{m+1}$. Since
\[
\frac{N_{m+1}}{N_m}\to1,
\]
we have $N_m\sim N$ and $\log N_m\sim\log N$. The prefix $w_N(\alpha)$ contains the prefix $w_{N_m}(\alpha)$, so
\[
p_{w_N(\alpha)}(k_{N_m})
\geq
(1-o(1))N_m
=
(1-o(1))N.
\]
Also
\[
k_{N_m}
=
(1+\varepsilon+o(1))\log_b N.
\]
Thus
\[
\delta(w_N(\alpha))
\geq
\frac{p_{w_N(\alpha)}(k_{N_m})}{k_{N_m}}
\geq
(1-o(1))\frac{N}{(1+\varepsilon)\log_b N}.
\]
Since $\varepsilon>0$ is arbitrary, taken over a countable sequence tending to $0$, we obtain
\[
\delta(w_N(\alpha))
\geq
(1-o(1))\frac{N}{\log_b N}
\]
almost surely. This proves the theorem. The lower bound for $\gamma$ follows from
\[
\delta(W)\leq\gamma(W).
\]
\end{proof}

\section{Summary of asymptotic scales}

The results above show that the online previous-copy complexity and the standard repetitiveness measures $\delta$ and $\gamma$ share the same natural scale in the worst case and almost surely:
\[
\frac{N}{\log_b N}.
\]

More precisely,
\[
\max_{|W|=N}\oc(W)
=
(1+o(1))\frac{N}{\log_b N},
\]
\[
\max_{|W|=N}\delta(W)
=
(1+o(1))\frac{N}{\log_b N},
\]
and
\[
\max_{|W|=N}\gamma(W)
=
(1+o(1))\frac{N}{\log_b N}.
\]
For Lebesgue-almost every $\alpha\in[0,1]$,
\[
\oc(w_N(\alpha))
=
\delta(w_N(\alpha))
=
\gamma(w_N(\alpha))
=
(1+o(1))\frac{N}{\log_b N}.
\]

The entropy spectrum and Hausdorff spectrum for $\oc$ are also determined:
\[
\lim_{N\to\infty}
\frac1N
\log_b
\#\left\{
W\in\{0,\ldots,b-1\}^N:
\oc(W)\leq
\kappa\frac{N}{\log_b N}
\right\}
=
\kappa
\]
for $0<\kappa\leq1$, and
\[
\dimH
\left\{
\alpha:
\liminf_{N\to\infty}
\frac{\oc(w_N(\alpha))\log_b N}{N}
\leq\kappa
\right\}
=
\kappa
\]
for $0\leq\kappa\leq1$.

\section*{Declaration of generative AI and AI-assisted technologies in the writing process}

During the preparation of this work, ChatGPT, by OpenAI, was used to assist with mathematical drafting, formalization, review, and editing. This work is shared as a preliminary AI-assisted mathematical note. The mathematical content may have been only partially reviewed and may contain errors; it should not be treated as peer-reviewed or as a fully verified manuscript.

\begin{thebibliography}{99}

\bibitem[KP18]{KP18}
D. Kempa and N. Prezza,
\emph{At the roots of dictionary compression: string attractors},
Proceedings of the 50th Annual ACM Symposium on Theory of Computing, 2018, 827--840.

\bibitem[KNP23]{KNP23}
T. Kociumaka, G. Navarro, and N. Prezza,
\emph{Toward a definitive compressibility measure for repetitive sequences},
IEEE Trans. Inform. Theory \textbf{69} (2023), no. 4, 2074--2092.

\bibitem[Pre17]{Pre17}
N. Prezza,
\emph{String attractors},
arXiv:1709.05314, 2017.

\bibitem[ZL77]{ZL77}
J. Ziv and A. Lempel,
\emph{A universal algorithm for sequential data compression},
IEEE Trans. Inform. Theory \textbf{23} (1977), no. 3, 337--343.

\end{thebibliography}

\end{document}
